% !TeX program = xelatex
\documentclass[zihao=-4, linespread=1.5]{ctexart}
\usepackage[top=2.5cm, bottom=2.5cm, left=3cm, right=2.5cm]{geometry}
\usepackage{amsmath, amssymb, bm}
\usepackage{graphicx}
\usepackage{fancyhdr}
\usepackage{titlesec}
\usepackage{cite}
\usepackage{hyperref}
\setlength{\parindent}{2em}
\setlength{\parskip}{0.5em}

% 页眉页脚
\pagestyle{fancy}
\fancyhf{}
\fancyhead[C]{\zihao{5} 学科技术前沿讲座课程大作业}
\fancyfoot[C]{\thepage}
\renewcommand{\headrulewidth}{0.4pt}

% 章节标题格式
\titleformat{\section}{\centering\bfseries\zihao{4}}{\thesection}{1em}{}
\titleformat{\subsection}{\bfseries\zihao{-4}}{\thesubsection}{1em}{}

\title{\vspace{2cm}\bfseries\zihao{3} 从计算机科学前沿探索到个人学术启航}
\author{王梓恒 \thanks{北京航空航天大学 计算机学院，北京 100191}}
\date{2026年5月}

\begin{document}

\maketitle
\vspace{1cm}

\begin{center}
\zihao{5}
\begin{tabular}{rl}
指导教师：& 李波 \\
所在院系：& 计算机学院 \\
专业名称：& 计算机科学与技术 \\
学\qquad 号：& 23371355 \\
\end{tabular}
\end{center}

\vspace{1cm}

\begin{abstract}
\noindent

本报告是修读《学科技术前沿讲座》课程后的综合总结。课程邀请了北京航空航天大学计算机学院和其他学院的多位学者，围绕计算机科学的核心问题、科研方法以及多个前沿技术领域进行深入讲解。内容涵盖计算机科学哲学与科研工作方法、虚拟现实与增强现实及数字孪生技术、人工智能与大模型、人形机器人与具身智能、高性能计算、多维度视觉计算以及存算传融合的自主计算架构等。

本文首先对各讲座的核心内容进行系统概括，梳理从学科通论到关键技术再到系统融合的知识脉络。然后重点阐述个人在学习过程中的深度思考与收获，包括对科研本质的重新认识、对人工智能热潮的理性分析、对技术范式转换的理解以及对个人未来学习和研究方向的启发。报告最后对课程进行总结，并展望智能化时代背景下计算机学子应有的定位与追求。

\vspace{0.5cm}
\noindent\textbf{关键词：} 学科前沿，人工智能，大模型，虚拟现实，具身智能，科研方法
\end{abstract}

% \newpage
% \clearpage
% \tableofcontents
% \clearpage
% \newpage

\section{绪论：讲座内容的多维概括}

\subsection{课程背景与总体框架}

学科技术前沿讲座课程旨在为计算机学院本科生提供了解学科最新发展动态、拓展学术视野、培养科研兴趣的平台。2026年春季学期，来自北航计算机学院及相关单位的八位专家学者围绕计算机科学的哲学基础、科研方法、人工智能（含大模型、具身智能）、虚拟现实、高性能计算、视觉计算、新型网络架构等主题做了系列报告。本报告是对这些讲座内容的系统总结与个人思考的集成。

\subsection{学科通论与科研方法论}

李波老师的开篇讲座《计算机科学的科研工作和论文写作》从哲学高度切入，探讨是否存在计算机科学这一根本问题，辨析广义计算机科学（与数学、物理、化学、社会科学交叉）与狭义计算机科学（聚焦计算机与计算本身）的矛盾。他指出当前我国计算机领域普遍存在重应用轻核心的现象，狭义计算机科学的研究人员过剩，而能够突破核心技术的人才匮乏。李老师认为科研的本质是“Search and Search Again”，创新点的提出是最大难点。在论文写作方面，他强调透彻的背景分析、翔实的数学模型以及有代表性的实验数据的重要性。这些观点为我们后续理解各具体领域的技术前沿提供了方法论指导。

\subsection{人工智能：从大模型到具身智能}

\subsubsection{人工智能发展状况与应用特点}

李波老师在另一场讲座《面向智能化发展的科技创新与人才培养》中系统梳理了人工智能从机械革命、电力革命、信息革命到当前智能化革命的演进脉络。他特别指出尽管人工智能潜力巨大，实际项目中人工智能的失败率高达百分之八十以上，主要原因在于数据问题、方法选择错误以及未能与具体业务场景深度结合。他提出二八现象，即通用大模型可能适合百分之八十的大众化应用，而百分之二十的复杂特殊问题需要专精特新的垂直领域模型。他还深入分析了人工智能面临的常识推理、不确定性推理等长期难题，以及深度学习作为静态数据拟合方法在安全攸关任务中的局限性，引导我们理性看待人工智能的能力与边界。

\subsubsection{大模型技术与DeepSeek案例}

李舟军教授的讲座《DeepSeek那些事儿》以现象级产品DeepSeek为案例，深入讲解大模型技术原理。他回顾了从专家系统到机器学习再到预训练加后训练范式的人工智能发展简史，重点剖析DeepSeek的技术创新，包括混合专家架构带来的训练与推理成本优势、GRPO强化学习方法赋予模型的强推理能力，以及全面开源策略带来的鲸鱼效应。李教授还结合自身团队的工作展示大模型在代码生成、社交网络内容理解、智能运维、网络威胁情报分析等垂直领域的落地应用新范式，让我们看到大模型从通用聊天走向行业专家的巨大潜力。

\subsubsection{具身智能与智能无人机}

刘偲老师的讲座《具身智能与智能无人机》将人工智能延伸到物理世界。她详细阐述具身智能的三大核心要素，包括智能体（由大语言模型或视觉语言模型驱动的大脑）、本体（人形机器人、无人机、机械臂等身体）以及数据与进化架构。她介绍国内外在具身智能领域的战略布局，并展示团队在无人机视觉语言动作模型、长距离规划导航以及多智能体协同等方面的前沿研究成果，生动诠释人工智能如何从虚拟大脑进化为能与物理世界交互的智能实体。

\subsection{虚拟现实与数字孪生}

周忠教授的讲座《虚拟现实技术前沿》系统介绍虚拟现实与增强现实技术的发展脉络、核心概念（沉浸感、交互性、构想性以及虚拟现实2.0提出的智能化、互通性、迭代性）与最新应用。他从太空算力、机器人谈到虚拟现实与增强现实在民生、工业、航天等领域的应用实例，深入介绍数字孪生技术及其在国防、地球科学、智能制造中的颠覆性作用。周教授还展示北航在虚拟现实专业课程建设与跨学科大作业方面的成果，体现虚拟现实作为解决问题思路的广阔前景。

\subsection{多维度视觉计算}

黄迪教授的讲座《多维度视觉计算》聚焦计算机视觉这一人工智能的核心领域。他从视觉计算的三个维度（目标、行为、事件）出发，探讨核心问题即视觉表征，介绍从生物启发的传统方法到基于预训练大模型的现代方法演进。黄教授团队在目标检测、识别、跟踪、分割、弱监督学习以及视觉内容生成等方面取得系列前沿成果，展现视觉计算从感知到生成再到理解的全面进步。

\subsection{高性能计算}

李波老师的讲座《从天河超级计算机看高性能计算技术》带我们回顾计算能力的基石。他以蛋白质结构预测、核爆模拟等大科学工程为引，阐述高性能计算对于国家战略与科技发展的重要性。讲座详细介绍天河一号、二号以及Blue Gene/L、Cray RedStorm、地球模拟器等典型机型架构，深入分析高性能计算的核心技术挑战，包括如何实现高生产率、高效率以及控制功耗。他探讨当前主流的集群体系结构、异构计算、高速互连网络、并行文件系统等技术趋势，让我们深刻理解支撑上层所有智能应用的算力是如何炼成的。

\subsection{存算传融合的自主计算}

张珊教授的讲座《存算传融合的自主计算》从网络互联的基本问题出发，回顾电路交换与分组交换的权衡，以及互联网的命名、寻址、路由、拥塞控制等经典难题。她敏锐指出当前传送算力存储三大资源体系独立演化的局限性，各自为政导致性能瓶颈。为此她提出存算传融合的创新理念，展示在芯片层面、架构层面以及系统层面的融合尝试。张教授描绘自主计算的愿景，即网络与设备融合后能够根据用户需求和环境状态自主按需安全地构建黑盒式计算环境。她还介绍团队在面向高时效视频分析的传算资源拟合、推理任务批处理机制等方面的具体研究成果。

\subsection{小结}

以上八位老师的讲座从哲学方法论到基础理论，从核心算法到系统架构，从软件智能到物理实体，共同勾勒出一幅计算机科学前沿发展的全景图。这些内容既相互独立又内在关联，为我后续深入思考个人收获与未来方向提供了丰富的素材。

\section{技术前沿深度分析：科研、AI与架构的范式思考}

\subsection{科研的真谛：从追随热点到追问本质}

李波老师的开篇讲座使我深刻认识到真正的科研不是简单学习最新工具或复现顶会论文，而是要对学科的根本问题提出追问。他提出的狭义与广义计算机科学的矛盾直击当前学术界的浮躁现象。多数计算机专业的研发人员不从事核心技术开发，非核心技术替代性强、价值低。这让我反思，我们日常所学的大量应用开发技能固然是工程实践的基础，但若止步于此便可能陷入万金油式的困境，缺乏核心竞争力。真正的科研应当敢于触碰那些难啃的骨头，即那些无法轻易被替代的核心技术。

在后续讲座中，这一观点成为我分析各项技术价值的照妖镜。当听到周忠教授展示虚拟现实在工业设计中的应用时，我不再仅仅惊叹于展示效果，而是思考背后的核心技术，包括高精度SLAM、物理仿真引擎、低延迟力反馈控制。同样，刘偲教授展示的具身智能机器人其核心在于大模型驱动的智能体、鲁棒的运动控制以及高效的数据进化架构。这些才是支撑起应用的基石。李老师强调科研工作需要理想主义的信念、持续不断的热情、完美主义的追求，这让我明白从事真正的科研意味着选择一条更艰难但更有价值的道路。我的核心收获在于，学会在未来的学习与研究中要有意识地去追问一项技术为什么能工作以及其根本瓶颈在哪里，而不是仅仅满足于怎么用它。这种从解决问题到发现问题的转变，才是科研思维的萌芽。

\subsection{人工智能热潮下的冷思考：能力、局限与价值}

多位老师并未一味鼓吹人工智能的万能，而是引导我们进行多维度的批判性思考。

能力与局限的辩证统一，李舟军教授清晰阐释大语言模型作为概率生成系统的本质，其智慧源于对海量数据的统计拟合而非真正的理解。这解释了大语言模型为何会产生幻觉，以及为何像9.8和9.11哪个大这样的简单问题会令其困惑。李波老师也指出深度学习模型是静态数据的拟合，难以适应动态变化，且不具备真正的逻辑推理和常识推理能力。这些分析让我认识到人工智能并非万能神药，其能力的边界恰恰源于其原理的局限性。

应用价值在于场景化，李波老师援引兰德公司报告指出人工智能项目高达百分之八十的失败率，核心原因往往不在于技术不够先进，而在于没有找准场景、数据质量差或与业务需求脱节。他强调智能具有场景化特征，人脸识别在闸机、抓拍机、巡飞弹上的实现细节天差地别。这让我明白真正的人工智能创新不是拿着锤子找钉子，而是深入理解具体场景的痛点，将技术与领域知识深度结合。ChatGPT的成功在于找到了聊天这个恰到好处的场景，而非强行用于高精度的专业诊断。

具身智能为人工智能带来新的可能性与挑战，刘偲老师的讲座展示具身智能如何通过与物理世界的交互为解决一些纯软件人工智能难以处理的问题提供新路径。但她也提到杨立昆的观点，当前机器人的常识不如一只猫。这揭示出具身智能在感知、决策、运动控制融合方面的巨大挑战，远未达到成熟。这场冷思考指导我如何评估一项人工智能技术的价值，不被华丽的演示迷惑，而是考察其数据依赖、泛化能力、鲁棒性以及在目标场景下的真实表现。

\subsection{技术范式的转换：从模块化到融合化}

对比不同老师的讲座，我观察到一种从模块化独立发展到融合化协同设计的深刻技术范式转换。

传送算存储的融合，张珊教授明确指出传统网络只负责传送，计算机只负责计算，存储系统只负责存储，这种独立演化的模式已无法满足未来应用对低延迟高效率的需求。她提出的存算传融合正是对这一问题的回应。无论是芯片级的存内计算、网络层的在网计算，还是系统级的天空计算，都是试图打破传统边界，通过协同设计来突破性能瓶颈。这不仅是技术的演进，更是一种系统思维的革新。

感知认知决策行动的融合，具身智能的出现正是感知、认知、决策与行动能力的深度融合。刘偲老师展示的AerialVLA模型让无人机通过统一的视觉语言动作模型同时完成环境理解、自然语言对话和飞行控制。这打破了传统机器人系统中感知规划控制串行管线的局限性。

软硬网的融合，高性能计算的发展也体现这种融合。李波老师介绍的异构计算以及天河、BlueGene等超级计算机的系统设计，无一不是算法、软件、芯片架构、互连网络和散热技术协同优化的结果。单一环节的突破难以带来系统级的性能飞跃。这一观察对我未来的学习具有重要方法论指导意义，在面对复杂系统问题时应当具备系统思维，不能局限于自己的小领域。解决方案的创新往往发生在多个层次的交叉地带。

\section{个人体会与收获：在巨变的时代中定位}

\subsection{知识体系的整合与重构}

通过本学期前沿讲座的学习，我最大的收获之一是原有的相对零散的计算机知识体系得到整合与重构。此前我学习操作系统、计算机网络、计算机组成原理、人工智能等课程时各门课相对独立。而本次讲座中，张珊教授从网络基本问题出发引出对存算传融合的需求，李波教授从高性能计算角度展示处理器、内存、网络互连如何协同，黄迪教授展示视觉算法如何在移动端和边缘端受限于算力和带宽，刘偲教授展示具身智能如何依赖高效的通信。这些内容将我脑海中的知识孤岛连接成一片完整的图景。我开始理解为什么操作系统中的虚拟内存管理要与网络协议栈中的RDMA结合，为什么计算机体系结构中的多核缓存一致性会影响分布式机器学习的效率，为什么计算机视觉中的模型压缩直接影响自主无人系统的响应速度。这种知识融合的视角将是我未来从事系统型研究的重要基础。

\subsection{科研思维与方法论的提升}

李波老师的科研方法论讲座对我影响最为深远。此前我常以为科研就是读论文加做实验加写论文的线性过程。李老师指出研究的核心在于搜索与比较，搜索所有可能的方法并比较，比较的结论引导新的发现。这让我意识到文献综述不仅仅是罗列前人工作，而是要通过深度比较发现现有方法的不足，从而找到自己的创新点。他还强调计算机科学研究的最难点是提出自己的创新点，而创新往往受限于基础科学的突破。这让我明白作为本科生虽然很难做出基础科学层面的突破，但可以在已有理论的基础上通过迁移、组合、改进现有方法在应用层面做出有价值的创新。关于论文写作，李老师强调数学模型往往是本质性的，这提醒我在今后的研究中要重视数学推导，而不仅仅满足于实验结果的漂亮。

\subsection{对人工智能未来发展的判断}

通过李舟军老师和刘偲老师的讲座，我对人工智能的未来发展有了更清醒的判断。通用大模型如DeepSeek-V3和GPT-4将作为基础设施存在，类似今天的操作系统或数据库，而真正的商业价值将爆发在通专结合的垂直领域，即那些需要深厚领域知识、高可靠性、实时性的场景如医疗诊断、工业控制、金融风控。这也意味着作为计算机专业的学生，如果仅仅掌握调用大模型应用程序接口的技能，未来很容易被替代。相反，如果能够深入某个垂直领域如生物信息、智能制造、航空航天，成为既懂人工智能又懂领域的复合型人才，将具备持久的竞争力。具身智能将是人工智能下一个十年的主战场，但其技术难度远超纯软件人工智能，需要硬件、算法、系统、通信等多方面的协同突破。我对刘偲老师团队在无人机视觉语言动作模型方面的工作非常感兴趣，希望未来有机会在该方向深入探索。

\subsection{对个人学习路径的规划}

结合以上收获，我对未来两年的学习与研究方向有了初步规划。
\begin{enumerate}
\item 夯实基础，继续深入学习操作系统、计算机组成原理、计算机网络、算法设计与分析等核心课程，尤其要理解底层原理，避免浮于表面。

\item 拓展交叉领域，选修或自学一些与人工智能加系统、人工智能加机器人相关的课程，如机器人学、控制理论、边缘计算等。同时关注生物信息、智慧城市等应用领域，培养跨学科思维。

\item 积极参与科研实践，争取进入实验室跟随老师或研究生参与实际项目。刘偲老师、黄迪老师、张珊老师等课题组的研究方向都非常前沿，我希望能够加入其中，从具体问题入手，锻炼发现问题和解决问题的能力。

\item 培养批判性思维，在阅读论文和聆听报告时有意识地追问这个方法的核心贡献是什么，它的局限在哪里，能否改进或迁移到其他场景。避免盲目崇拜顶会论文或工业界巨头的工作。

\item 关注国产技术生态，李舟军老师提到DeepSeek与华为昇腾的合作，张珊老师提到自主计算对国产芯片的意义。我认识到在当前国际环境下掌握自主可控的关键技术至关重要。我会关注国产操作系统如OpenHarmony、国产人工智能芯片如昇腾和寒武纪以及国产数据库和中间件，为将来投身国家信息技术创新事业打下基础。

\end{enumerate}
\subsection{课程带来的职业与人生启示}

这次课程也让我对未来的职业选择有了更深的思考。李波老师提到核心技术和非核心技术的差异在于不可替代性。虽然短期内从事应用开发可能更容易找到工作，但长期来看掌握核心技术的人才才能真正立于不败之地。周忠教授展示的虚拟现实与增强现实在医疗、航天等领域的应用，刘偲教授展示的无人机在应急救援、物流配送中的价值，让我感受到技术从业者的社会责任感。我们不仅仅是在写代码和调参数，更是在用技术解决实际问题、改善人们的生活、甚至保卫国家安全。这种使命感将激励我在未来的学习和工作中始终牢记将论文写在祖国的大地上。

\section{结论}

通过本学期学科技术前沿讲座课程的学习，我系统接触了计算机科学多个前沿领域的最新进展，包括科研方法论、人工智能（大模型与具身智能）、虚拟现实与数字孪生、多维度视觉计算、高性能计算以及存算传融合的自主计算架构。各位老师的精彩报告不仅传递了知识，更启迪了思维。

我深刻认识到计算机科学是一门既需要深邃哲学思考又需要扎实工程实践的学科。在人工智能浪潮席卷全球的今天，我们既要积极拥抱新技术带来的机遇，也要保持清醒的头脑，理性分析技术的边界与局限。真正的创新往往发生在不同学科的交叉点上，发生在对现有技术范式的反思与突破中。作为新时代的计算机学子，我们应当树立远大理想，勇于挑战核心技术难题，将个人发展融入国家科技自立自强的伟大征程。

本次课程的学习是我本科阶段一次宝贵的思想洗礼。我将以此次学习为起点，在未来的日子里继续深入钻研专业知识，培养批判性思维和创新精神，努力成为一名既有扎实理论基础又有解决实际问题能力的计算机科学工作者。

\end{document}